\section{Related Work}
\label{sec:related}

\para{Linear probes and emergent representations.} \citet{alain2016linear} established the linear-probe methodology: train a linear classifier on frozen activations at each layer to test which information is linearly decodable. \citet{li2022emergent} showed the strongest form of the latent-variable claim on Othello-GPT, where a transformer trained only on move sequences develops a probe-decodable and causally interveneable board state. Our work uses linear probes in the same spirit, but to test whether multiple modes of \emph{text production} (not a single task variable) coexist as disentangled directions in the same residual stream.

\para{Stylometry and authorship attribution.} A long literature treats authorial style as a latent variable recoverable from surface text. \citet{rivera2021universal} learn universal authorship embeddings that transfer across domains (reviews, fanfiction, Reddit). Most relevant to us, \citet{sarfati2025prompt} show that literary style in 13 novels is linearly decodable from \llama last-token embeddings, reaches its sharpest separation at intermediate layers, and lies within $\sim$16 PCA dimensions of class-mean differences. We extend this template along two axes: we test \emph{multiple} modes within the same embedding cache, and we measure their geometric \emph{relationships} (orthogonality and transfer), which a single-mode study cannot.

\para{AI-text detection.} A fast-moving subfield treats AI-vs-human classification as its own problem. \citet{mitchell2023detectgpt} introduce DetectGPT, a zero-shot detector based on local curvature of the LM's output distribution. \citet{verma2023ghostbuster} and others fit supervised classifiers on hand-crafted or neural features. \citet{schafer2025stylistic} obtain F1 above 0.90 with 220 hand-crafted stylistic features, noting that AI text systematically biases towards greater lexical complexity and lower variation. None of these methods asks whether the AI-vs-human axis shares geometry with stylistic or surface-form axes. Our cross-mode transfer experiment is, to our knowledge, the first to quantify the overlap between an AI-text direction and a register direction within the same model.

\para{Feature discovery via sparse autoencoders.} \citet{templeton2024scaling} train sparse autoencoders on Claude 3 Sonnet and recover interpretable features for cities, emotions, code, and deception. SAEs are an unsupervised counterpart to probing: they discover candidate latent variables without labels. Our supervised probing study complements this line by directly testing whether labeled \modes are linearly decodable and whether their directions are mutually orthogonal --- a property that SAE dictionaries are often assumed but rarely verified to satisfy across feature pairs.

\para{Datasets and corpora.} Authorship work commonly uses Project Gutenberg \citep{sarfati2025prompt}. AI-text work commonly uses HC3 \citep{guo2023chatgpt}, a 24K Q\&A corpus of paired human and ChatGPT answers across five domains. We use both, and add two synthetic corpora --- one of formal-vs-casual essay pairs generated by \gptmini for register control, and one of clean text with QWERTY-adjacent versus Dvorak-adjacent typo injection. Synthetic data lets us isolate a single axis of variation; the cost is that the resulting \modes reflect our construction rather than naturally-occurring writing, a tension we discuss in \secref{sec:discussion}.

\para{Position of our work.} Prior work establishes that individual production-related axes (style, authorship, AI-vs-human) are linearly recoverable from hidden states. We show that these are best understood as instances of a unified \modes abstraction: distinct \modes coexist in the residual stream as near-orthogonal directions, and the partial overlaps that do exist (notably \register $\to$ \aihuman) trace to identifiable content correlations rather than to a shared underlying axis.
